\phantomsection
\section*{十一、总结}
\addcontentsline{toc}{section}{十一、总结}

本文围绕业界在资源自适应模型切分算法、业务感知 batching 优化算法、用户移动场景下请求调度算法三个方向的学术研究和工程应用进行了综述。与传统数据中心 LLM serving 相比，CEC 场景需要额外考虑端-边-云异构资源、无线链路波动、边缘节点负载变化、用户移动、KV/状态连续性、隐私边界和业务 QoS。现有研究已经分别在异构推理、LLM serving runtime、KV/cache 管理和移动边缘服务连续性方面提供了可借鉴基础，但仍需要结合无线场景做系统级集成 \refciterange{3}{7}, \refcite{13}, \refcite{18}, \refciterange{21}{26}。


总体来看，三类算法的落点是不同时间尺度上的协同控制：模型切分解决“在哪里跑”，batching 解决“怎么排”，请求调度解决“由谁来跑”。如果把这三个问题分开优化，容易得到局部最优；如果把它们放入统一的分层闭环中，则更容易得到可部署、可解释、可回退的工程方案。


面向华为无线软件实现，推荐采用与 12 个月项目周期对齐的落地路线：T 到 T+2 月完成技术调研、方案选择、验证环境梳理和最小 serving baseline；T+3 到 T+8 月并行推进资源自适应模型切分算法与业务感知 batching 算法，形成可运行源码、必要的推理框架适配、测试结果和分析报告；T+9 到 T+12 月在前两项算法形成的实例画像、部署模板和 batching 行为基础上，完成用户移动场景下的请求调度算法。该路线能够避免从零构建推理系统，同时保证三项算法都对应清晰的基线、指标、源码交付和验证场景。




